% This Source Code Form is subject to the terms of the Mozilla Public
% License, v. 2.0. If a copy of the MPL was not distributed with this
% file, You can obtain one at https://mozilla.org/MPL/2.0/.
\documentclass[11pt,a4paper]{article}

\usepackage[margin=2.5cm]{geometry}
\usepackage{amsmath,amssymb,amsthm}
\usepackage{booktabs}
\usepackage{array}
\usepackage{microtype}
\usepackage[hidelinks]{hyperref}

\emergencystretch=3em

\newcommand{\lean}[1]{\texttt{#1}}
\newcommand{\R}{\mathbb{R}}
\newcommand{\N}{\mathbb{N}}
\newcommand{\Prob}{\mathbb{P}}
\newcommand{\none}{\bot}

\theoremstyle{definition}
\newtheorem{definition}{Definition}[section]
\newtheorem{theorem}[definition]{Theorem}
\newtheorem{remark}[definition]{Remark}

\title{The problem that unisono solves}
\author{}
\date{}

\begin{document}
\maketitle

This document specifies the problem that unisono solves, independently of
any solution. The Lean files in \lean{lean/Unisono/Spec} state the same
definitions; each definition below names its Lean counterpart. Internal
mechanisms such as schedules, clock estimators, controllers and routing are
left to the solution and do not appear here. Statements about the problem,
such as limits that no solution can beat, and statements about a solution,
such as what it achieves, are made against this specification and under a
model of the world (Section~\ref{sec:models}). A solution is one whole
system, judged by what it achieves for the objectives of
Section~\ref{sec:objectives}.

\paragraph{The objectives at a glance.} Each objective accrues a cost per
frame of the source or per output sample of a speaker, and is measured as a
long-run average per second of true time. Section~\ref{sec:objectives}
defines each symbol where it is used: $R$ is the number of speakers,
$T_r(n)$ the time at which speaker $r$ plays frame $n$, $\ell_r(n)$ its
latency, $L^*(n)$ the target latency, $v_r(k)$ the playback speed at output
sample $k$, $j$ the next step of continuous playback after $k$, and
$\Delta t$ the time between the two steps. The last two columns give the
weight of each objective in each mode.

\begin{center}
\footnotesize
\begin{tabular}{@{}l>{$}l<{$}ll>{$}c<{$}>{$}c<{$}@{}}
\toprule
Objective & \text{cost accrued} & unit & Lean & \text{low latency} & \text{lip sync} \\
\midrule
Synchronization & \frac{1}{R^2}\sum_{r<r'} (T_r(n)-T_{r'}(n))^2 & s$^2$ & \lean{syncCost} & w_S & w_S \\[2pt]
Latency & \frac{1}{R}\sum_r \ell_r(n) & s & \lean{latencyCost} & w_\Lambda & 0 \\[2pt]
Tracking & \frac{1}{R}\sum_r (\ell_r(n)-L^*(n))^2 & s$^2$ & \lean{trackingCost} & 0 & w_\Theta \\[2pt]
Roughness & (v_r(j)-v_r(k))^2/\Delta t & s$^{-2}$ & \lean{roughnessCost} & w_V & w_V \\
Missed samples & \text{samples not played as captured} & s$^{-1}$ & \lean{missedAt} & w_M & w_M \\
Discontinuities & \text{breaks in playback} & s$^{-1}$ & \lean{discontinuityAt} & w_D & w_D \\
\bottomrule
\end{tabular}
\end{center}

\section{Scope}

One sender captures audio from one source. Several speakers play it. Each
speaker is a playback device driven by a host; a host may drive several
speakers, and the sender's host may drive speakers too. Every speaker should
play every sample of the source
\begin{itemize}
\item synchronized: all speakers play a sample at the same time;
\item with low latency: soon after the sample was captured, or, in lip-sync
  mode, at a target latency that the application sets;
\item with little rate distortion: at the speed at which the source was
  captured, with the speed changing only slowly;
\item without missed samples and without breaks: every output sample plays
  the original sound in its place, continuously.
\end{itemize}
Out of scope: correcting for the time that sound takes from a speaker to a
listener (room correction); what a speaker plays instead of a missed sample
(concealment), since every output sample that is not the original counts as
missed whatever it plays; the quality of resampling, apart from its
lookahead (Section~\ref{sec:solutions}); and devices of different nominal
rates.

\section{The world}\label{sec:world}

The world is everything that no solution chooses (\lean{World}). A model of
the world (Section~\ref{sec:models}) says which worlds can happen and with
which probabilities.

\begin{definition}[World]
\begin{itemize}
\item \emph{True time} $t \in \R$ is measured in seconds (\lean{Time}). No
  device can read it.
\item \emph{The capture device} of the sender (\lean{CaptureDevice})
  captures frame $n \in \N$ of the source at true time $c(n)$, strictly
  increasing in $n$ (\lean{capture}). A frame holds one sample of each
  channel. The device hands the frames to the sender's host in blocks:
  block $b$ holds the frames $f(b) \le n < f(b+1)$ (\lean{first}) and is
  handed over at true time $h(b) \ge c(f(b+1)-1)$ (\lean{handed}). For each
  block the device reports a timestamp $\hat c(b)$, a reading of its host's
  clock that estimates when the block's first frame was captured
  (\lean{reported}); it may be wrong.
\item \emph{Speaker} $r \in \{0,\dots,R-1\}$ is a playback device
  (\lean{PlaybackDevice}). Its output sample $k \in \N$ leaves the device at
  true time $o_r(k)$, strictly increasing in $k$ (\lean{out}). The device
  requests block $b$, the output samples $g_r(b) \le k < g_r(b+1)$
  (\lean{first}), at true time $a_r(b) < o_r(g_r(b))$ (\lean{asked}), and
  reports a timestamp $\hat o_r(b)$ that estimates when the block's first
  output sample leaves (\lean{reported}).
\item \emph{Every host} $i$ has a clock $C_i$ (\lean{Clock}), a continuous,
  strictly increasing function from true time to readings.
\item \emph{The network} (\lean{Network}) carries messages between any two
  hosts. When host $i$ sends its $q$-th message to host $j$ at true time
  $s$, the network either delivers it at a true time
  $\mathrm{deliver}_{ij}(s,q) > s$ or drops it (\lean{deliver}). Delivery
  and loss may depend on $s$ and $q$, but not on the payload.
\item \emph{The application} sets a target latency $L^*(n)$ for each frame
  $n$, used in lip-sync mode, and gives it to the sender's host together
  with the frame (\lean{target}).
\item All devices have the same nominal rate $\rho$, in frames per second
  (\lean{rate}); each runs at the rate of its own oscillator. The hosts are
  the sender's host and the speakers' hosts (\lean{hosts}).
\end{itemize}
\end{definition}

\begin{definition}[Positions in the source]
Speakers resample: an output sample plays the source at a position
$x \ge 0$, counted in frames, which in general is not a whole number. The
sound at position $x$ counts as captured at
\[
  c(x) = c(\lfloor x \rfloor) + (x - \lfloor x \rfloor)\,
         \bigl(c(\lfloor x \rfloor + 1) - c(\lfloor x \rfloor)\bigr),
\]
where $\lfloor x \rfloor$ is $x$ rounded down (\lean{World.captured}).
Position $999.98$ lies $98\%$ of the way from frame $999$ to frame $1000$,
and counts as captured $98\%$ of the way between their capture times.
\end{definition}

\section{What a host observes}

A host acts at \emph{events} (\lean{Observation}, \lean{Event}). At each
event it observes one of the following, together with the reading of its
own clock:
\begin{itemize}
\item the capture device hands over a block: its frames, each with its
  target latency, and the block's timestamp (at the sender's host only);
\item a speaker's device requests a block: its output samples and the
  block's timestamp;
\item a message arrives: its sender, its payload, and the frames of the
  source whose data it carries;
\item a timer that the host set fires.
\end{itemize}
The \emph{view} of a host at an event is the list of all events it has
observed so far, each with its clock reading. True time is not in the view,
and neither is anything that happens at another host unless a message
reports it.

\section{Solutions and their run}\label{sec:solutions}

\begin{definition}[Solution]
At each event, a host sends messages to any hosts, with a payload of the
solution's choice and with the data of frames that it holds; sets timers at
readings of its clock; and, when a speaker's device requests a block,
chooses for each output sample of the block what it plays: the source at a
position $x$, or anything else, written $\none$ (\lean{Action}). A
\emph{solution} $\sigma$ determines what each host does at each event as a
function of its view (\lean{Solution}). A host \emph{holds} a frame once an
event has brought its data: the capture device's handoff at the sender, or
a message.
\end{definition}

Four limits apply to every solution.
\begin{enumerate}
\item \emph{Causality.} What a host does at an event depends only on its
  view.
\item \emph{Data.} A host sends only frames it holds. A speaker's resampler
  computes the sound at position $x$ from the frames $n$ with $|n - x| < a$,
  where the integer $a \ge 1$ is the resampler's lookahead: $a = 1$ for
  linear interpolation (position $999.98$ needs frames $999$ and $1000$),
  about half the filter length for a longer filter. An output sample plays
  the original sound only if its host holds all of those frames when the
  device requests the block; otherwise it is a missed sample
  (\lean{playable}). The lookahead $a$ is a parameter.
\item \emph{Capacity.} In any second, each host sends each other host at
  most $M$ messages carrying at most $F$ frames in total; $M$ and $F$ are
  parameters (\lean{Solution.WithinBudget}).
\item \emph{Speed bound.} At every step of continuous playback, every
  speaker's speed $v$ (Section~\ref{sec:objectives}) satisfies
  $|v - 1| \le u_{\max}$, where the parameter $u_{\max}$ is the
  resampler's limit (\lean{Solution.WithinSpeed}).
\end{enumerate}

\begin{definition}[Run]
The \emph{run} of a solution in a world (\lean{Solution.states}) processes
events in chronological order. At each step the earliest pending event is
handled: the capture device's next block, a speaker's device's next
request, the arrival of a message in transit, or a timer that is set. The
world determines when each occurs: the devices determine when they hand
over and request blocks, the network determines when a message arrives or
that it is dropped, and a host's clock determines when it reaches the
reading of a timer. The host adds the event to its view and does what the
solution chooses for that view. Simultaneous events are processed in a
fixed order: the capture device's first, then the speakers' devices' in the
order of the speakers, then messages and timers in the order in which they
were sent or set. What speaker $r$ plays at output sample $k$ is what was
decided when its device requested the block holding $k$
(\lean{Solution.plays}).
\end{definition}

Every solution therefore has exactly one run in every world.

\begin{theorem}\label{thm:run}
The run processes events in chronological order
(\lean{Solution.eventTime\_le\_succ}). If the run passes every time, that is,
its events eventually occur later than any given true time, then the
solution decides what every output sample of every speaker plays
(\lean{Solution.decided}).
\end{theorem}

A solution that schedules infinitely many events before some finite time
$t$, for example timers ever closer together, never passes $t$, and every
output sample after $t$ counts as missed.

\section{The objectives}\label{sec:objectives}

All objectives are measured in true time from what each speaker plays. Let
$p_r(k) \in \R \cup \{\none\}$ be what speaker $r$ plays at output sample
$k$: a position in the source, or $\none$ for anything else.

\begin{definition}[Continuous playback]
Playback of speaker $r$ is continuous from output sample $k$ to $k+1$ when
both play the source and the position moves strictly forward by less than
two frames, $0 < p_r(k+1) - p_r(k) < 2$ (\lean{PlaysContinuously}): no frame
is skipped and none is played twice.
\end{definition}

\begin{definition}[Play time and latency]
Speaker $r$ plays frame $n$ at
\[
  T_r(n) = o_r(k) + \frac{n - p_r(k)}{p_r(k+1) - p_r(k)}\,\bigl(o_r(k+1) - o_r(k)\bigr),
\]
where $k$ is the first continuous step with $p_r(k) \le n \le p_r(k+1)$
(\lean{World.playTime}); if there is no such step, speaker $r$ does not play
frame $n$. The latency of frame $n$ at speaker $r$ is
$\ell_r(n) = T_r(n) - c(n)$ (\lean{World.latency}).
\end{definition}

\begin{definition}[Synchronization]
For frame $n$,
\[
  S(n) = \frac{1}{R^2} \sum_{r < r'} \bigl(T_r(n) - T_{r'}(n)\bigr)^2,
\]
summed over the pairs of speakers that both play the frame
(\lean{World.syncCost}). When every speaker plays it, $S(n)$ is the
variance of the times $T_r(n)$ (\lean{sum\_pairs\_eq}). It compares speakers
only with each other.
\end{definition}

\begin{definition}[Latency and tracking]
For frame $n$,
\[
  \Lambda(n) = \frac{1}{R}\sum_r \ell_r(n),\qquad
  \Theta(n) = \frac{1}{R}\sum_r \bigl(\ell_r(n) - L^*(n)\bigr)^2,
\]
each summed over the speakers that play the frame (\lean{World.latencyCost},
\lean{World.trackingCost}). A speaker that does not play the frame adds
nothing to $S$, $\Lambda$ or $\Theta$; it counts missed samples instead.
\end{definition}

\begin{definition}[Speed]
On a continuous step $k$ of speaker $r$, the speed is
\[
  v_r(k) = \frac{c(p_r(k+1)) - c(p_r(k))}{o_r(k+1) - o_r(k)},
\]
the capture time that the step plays per second of playback
(\lean{World.speed}). A speed of one plays the source at the pitch and
tempo at which it was captured; a speed $v$ shifts the pitch by
$1200\log_2 v$ cents. The speed is one minus the rate of change of the
latency, $v = 1 - \dot\ell$ (\lean{speed\_eq}).
\end{definition}

\begin{definition}[Roughness]
For a continuous step $k$ of speaker $r$, let $j > k$ be the next
continuous step. The roughness between them is
\[
  V_r(k) = \frac{\bigl(v_r(j) - v_r(k)\bigr)^2}{\Delta t},\qquad
  \Delta t = \frac{o_r(j) + o_r(j+1)}{2} - \frac{o_r(k) + o_r(k+1)}{2},
\]
where $\Delta t$ is the time between the midpoints of the two steps
(\lean{World.roughnessCost}). In continuous playback $j = k + 1$, and the
sum of $V_r$ is the integral $\int \dot v^2\,dt$. Across a break, $V_r(k)$ is
what a speed changing linearly over the break would cost, which is the least
that any change of speed over that time costs; a break therefore does not
hide a change of speed (\lean{roughnessCost\_across}).
\end{definition}

\begin{definition}[Missed samples]
At output sample $k$ of a speaker (\lean{missedAt}):
\begin{itemize}
\item an output sample that plays $\none$ counts one;
\item an output sample that plays position $y$, when the last earlier output
  sample that played the source played position $x$, $g$ output samples
  before it, counts the frames that $y$ skips beyond those $g$, namely
  $\max(\lfloor y - x - 1 \rfloor - g, 0)$, plus the frames it plays again,
  $\max(\lfloor 1 - (y - x) \rfloor, 0)$.
\end{itemize}
\end{definition}

\begin{center}
\small
\begin{tabular}{@{}p{0.55\textwidth}cl@{}}
\toprule
What a speaker does & Missed samples & Lean \\
\midrule
plays continuously & $0$ & \lean{missedAt\_continues} \\
jumps forward over $J$ frames & $J$ & \lean{missedAt\_skip} \\
plays $J$ frames again & $J$ & \lean{missedAt\_repeat} \\
plays $g$ samples of anything else, then continues where it would have been & $g$ & \lean{missedAt\_after\_gap} \\
plays $g$ samples of anything else, then resumes where it stopped & $g$ & \lean{missedAt\_resume} \\
\bottomrule
\end{tabular}
\end{center}

A speaker that resumes where it stopped plays everything after a little
later; that costs synchronization and latency, not missed samples.

\begin{definition}[Discontinuities]
Playback of speaker $r$ breaks after output sample $k$ when $p_r(k)$ plays
the source and the step from $k$ to $k+1$ is not continuous: the next output
sample plays $\none$, jumps forward by two frames or more, or jumps back
(\lean{discontinuityAt}). Each break counts one discontinuity. A gap counts
one however long it lasts (\lean{discontinuities\_gap}), and a jump counts
one (\lean{discontinuityAt\_jump}).
\end{definition}

Missed samples measure how much sound is lost; discontinuities measure how
often playback breaks, each break being an audible click.

\begin{definition}[Long-run averages and cost]
A cost of each frame accrues until true time $T$ as
\[
  A(T) = \sum_{n:\,c(n) < T} \mathrm{cost}(n)\,\bigl(c(n+1) - c(n)\bigr),
\]
an integral over the time of capture (\lean{World.accruedFrames}). A cost
of each output sample accrues as
\[
  A(T) = \frac{1}{R}\sum_r \sum_{k:\,o_r(k) < T} \mathrm{cost}_r(k)
\]
(\lean{World.accruedOutputs}). The long-run average of $A$ is
$\limsup_{T \to \infty} A(T)/T$ (\lean{longRun}). The cost of a playback,
for weights $w = (w_S, w_\Lambda, w_\Theta, w_V, w_M, w_D) \ge 0$
(\lean{Weights}), is
\begin{multline*}
  J = \limsup_{T \to \infty} \frac{1}{T}\Bigl(
      w_S A_S(T) + w_\Lambda A_\Lambda(T) + w_\Theta A_\Theta(T) \\
    + w_V A_V(T) + w_M A_M(T) + w_D A_D(T) \Bigr),
\end{multline*}
where $A_S$, $A_\Lambda$ and $A_\Theta$ accrue $S$, $\Lambda$ and $\Theta$
per frame, and $A_V$, $A_M$ and $A_D$ accrue roughness, missed samples and
discontinuities per output sample (\lean{World.cost}). The cost of a
solution in a world is the cost of what its speakers play in its run there
(\lean{Solution.cost}); when the world is random, a solution's cost is its
expected cost (\lean{Solution.expectedCost}).
\end{definition}

The weights are parameters: a solution is designed for weights, and
statements about solutions say for which. Low-latency mode has
$w_\Theta = 0$; lip-sync mode has $w_\Lambda = 0$. How the objectives
combine may be revisited once the limits and the solutions are known.

\begin{remark}[Muting a speaker]
A speaker that does not play a frame adds nothing to that frame's
synchronization, latency and tracking costs. Not playing frame $n$ at
speaker $r$ lowers $\Lambda(n)$ by exactly $\ell_r(n)/R$
(\lean{latencyCost\_unplayed}). When all other speakers play the frame at
one time and speaker $r$ is off from them by $\Delta$, not playing it at
speaker $r$ lowers $S(n)$ by $(R-1)\Delta^2/R^2$
(\lean{syncCost\_one\_off}). Muting speaker $r$ for a while costs one missed
sample per output sample, so it lowers the cost only if
\[
  w_M\,\rho < w_\Lambda\,\ell + w_\Theta\,(\ell - L^*)^2 + w_S\,\frac{(R-1)\,\Delta^2}{R},
\]
where $\rho$ is the sample rate, $\ell$ the speaker's latency and $\Delta$
its offset from the others. At $\rho = 48\,000$ per second and
$\ell = 50$\,ms, with no tracking or synchronization error, this needs
$w_M$ below a millionth of $w_\Lambda$. Weights are chosen so that it does
not hold.
\end{remark}

\begin{remark}[Dilution of one speaker]
With the normalization by $R^2$, one speaker off by $\Delta$ from all
others adds $(R-1)\Delta^2/R^2 \approx \Delta^2/R$ to the synchronization
cost: the variance weighs a single outlying speaker less as the number of
speakers grows. Synchronization is measured as this variance, over all
pairs.
\end{remark}

\section{Rate distortion}

The earlier work that this project started from charges the \emph{warp}
$u$, the fraction by which a receiver's resampler plays faster or slower
than its nominal rate, with $\int \dot u(t)^2\,dt$, the $\dot H^1$
seminorm. It charges the derivative and not the warp itself: fast changes
of pitch are heard, a slowly varying small offset is not. Its bound on the
warp, half a percent, is the resampler's limit; here it is the speed bound
$u_{\max}$ of Section~\ref{sec:solutions}.

Here the speed is that of the source as heard, against its capture in true
time; its steady value is one, and the clocks' drift appears only in what a
solution does to keep it at one. The earlier warp is the resampler's rate
against its own device, whose steady value is the drift. The two
derivatives differ only by how fast the clocks' own rates change, so the
roughness here is the earlier warp cost, measured on what is heard.

Without the speed bound, the best recovery from an excess latency
$\Delta L$ under latency and roughness costs alone plays off pitch by an
amount that grows as $\Delta L^{3/4}$, without limit. The bound caps it, as
in the earlier work. Open for discussion: whether the weighting of a speed
modulation by the square of its rate matches how listeners detect frequency
modulation, or a weighting by modulation rate should replace it.

\section{Models of the world}\label{sec:models}

A model of the world says which worlds can happen and, when the world is
random, with which probabilities. Statements hold under a model, and the
same solution can do differently under different models. Each model is a
file of its own in \lean{lean/Unisono/Spec/Models}, and each comes with an
example showing that it is not empty.

\subsection{Bounded worlds}

The model \lean{Bounded} is deterministic: a statement under it holds for
every world within the bounds (\lean{World.WithinBounds}), the worst case.
Its parameters (\lean{Bounds}) are the rate errors
$\varepsilon_C, \varepsilon_D \in [0, 1)$, the timestamp error
$\varepsilon_R \ge 0$, the handoff delay $\eta \ge 0$, the request leads
$0 < \lambda_{\min} \le \lambda_{\max}$, the block size $m_B \ge 1$, the
delays $0 < d_{\min}(i,j) \le d_{\max}(i,j)$ of each link, and the number
$m_L$ of consecutive losses. A world is within the bounds when:
\begin{itemize}
\item clocks: $|C_i(t) - C_i(s) - (t - s)| \le \varepsilon_C\,(t - s)$ for
  $s \le t$;
\item device rates: $|c(n+1) - c(n) - 1/\rho| \le \varepsilon_D/\rho$, and
  likewise for $o_r$;
\item timestamps: $|\hat c(b) - C_{\mathrm{sender}}(c(f(b)))| \le
  \varepsilon_R$ and $|\hat o_r(b) - C_{\mathrm{host}(r)}(o_r(g_r(b)))| \le
  \varepsilon_R$;
\item handoff: $h(b) - c(f(b+1) - 1) \le \eta$;
\item request lead: $\lambda_{\min} \le o_r(g_r(b)) - a_r(b) \le
  \lambda_{\max}$;
\item block size: $f(b+1) - f(b) \le m_B$ and $g_r(b+1) - g_r(b) \le m_B$;
\item delays: every delivered message from $i$ to $j$ takes between
  $d_{\min}(i,j)$ and $d_{\max}(i,j)$;
\item losses: of any $m_L + 1$ consecutive messages on a link, at least one
  is delivered, whenever they are sent.
\end{itemize}
A random world can be required to be within bounds almost surely
(\lean{AlmostSurelyWithin}). The model is not empty
(\lean{Bounds.example\_within}).

\subsection{Delays, pauses and losses}

The model \lean{Delays} gives the network a probability distribution; it is
the model of the earlier work. Each link from host $i$ to host $j$
(\lean{Link}) has a baseline $B(u) > 0$, a jitter $J(u) \ge 0$, a set of
episodes, each a pause or an outage with a start and a length $H > 0$, and a
set of individually dropped messages.

\begin{definition}[Delivery]
A message sent at true time $s$ is released at
\[
  U(s) = \inf\{u \ge s : \text{no pause is under way at } u\}
\]
(\lean{Link.release}). The $q$-th message on the link is dropped if it is
one of the individually dropped messages or if an outage is under way at
$U(s)$; otherwise it arrives at $U(s) + B(U(s)) + J(U(s))$
(\lean{Link.deliver}).
\end{definition}

A pause thus holds back every message sent during it and releases them all
when it ends: the first message sent in a pause of length $H$ is late by
$H$, and messages released together arrive together, the earlier work's
zero-drain idealization. Two messages sent on the same link at the same
instant take the same time.

\begin{definition}[The model's distribution]
The parameters (\lean{DelayModel}) are the slope $\beta \ge 0$ of the
baseline, the tail constants $C > 0$ and $x_0 > 0$ of the jitter, the rate
$\lambda_\kappa \ge 0$ and the length distribution $F_\kappa$ on
$(0, \infty)$ of each kind $\kappa$ of episode, and the loss probability
$p \in [0, 1]$. A random network follows the model
(\lean{DelayModel.Holds}) when:
\begin{itemize}
\item the baseline changes slowly: $|B(u) - B(u')| \le \beta\,|u - u'|$;
\item the jitter comes arbitrarily close to zero, $\Prob(J(u) < \delta) > 0$
  for every $\delta > 0$, so the baseline is the least time a message takes;
\item large jitter is rare: $\Prob(J(u) > x) \le C e^{-x/x_0}$ for $x \ge 0$;
\item episodes of kind $\kappa$ start at rate $\lambda_\kappa$ with lengths
  distributed as $F_\kappa$: the expected number that start within any $t$
  seconds with length in a set $A$ is $\lambda_\kappa\, t\, F_\kappa(A)$;
\item each message is dropped individually with probability $p$;
\item the state of all links over time, $X(u)$, made of each link's jitter
  and whether a pause or an outage is under way, is jointly
  \emph{stationary}: the distribution of the path $(X(u + \tau))_u$ equals
  that of $(X(u))_u$ for every $\tau$; and \emph{mixing}: for any two
  events $E$ and $E'$ about the path,
  $\Prob(X \in E,\ \theta_\tau X \in E') \to \Prob(X \in E)\,\Prob(X \in
  E')$ as $\tau \to \infty$, where $\theta_\tau$ shifts the path by $\tau$
  (\lean{Stationary}, \lean{Mixing});
\item the individually dropped messages of all links are jointly stationary
  and mixing in the message number.
\end{itemize}
A random world follows the model when its network consists of such links
(\lean{DelayModel.Random}).
\end{definition}

Mixing implies that long-run averages equal their expectations almost
surely. This rules out worlds that, with some probability, drop every
message or hold every message forever. Nothing else is assumed independent:
links may be alike, episodes may come in clusters, and the jitter, the
episodes and the losses may depend on each other. Mixing is assumed without
a rate; a rate will be added when a statement about finite time windows
needs one. The model is not empty (\lean{DelayModel.example\_holds}).

Models of how the clocks wander come with the statements that need them.

\end{document}
